\chapter{Производная относительная инволюция и общая норма кривизны}
\label{ch:v7-derived-relative-involution-curvature-norm-gate}

\section{Положительная замена ролевого суперследа}

Пусть $A:\mathbb C^{11}\to\mathbb C^{10}$ --- ориентированная половина
физического incidence-оператора, а $A_0$ --- единичный фон семи корневых
рёбер, уже выбранных старым $H_{15}$ и бесцветным проектором. Определим
\begin{equation}
 \mathcal S_V(A)=\frac12\left(
 \|A^*A-A_0^*A_0\|_{\rm HS}^2+
 \|AA^*-A_0A_0^*\|_{\rm HS}^2\right).
 \label{eq:v7-relative-norm-vertex-action}
\end{equation}
Это положительная норма относительной Hodge-кривизны без суперследа.

Рёберный блок использует уже выведенную инволюцию и полные Casimir-веса:
\begin{equation}
 \mathcal S_E=
 \sum_{a=1}^7m_a(|z_a|^2-1)^2
 +\sum_{j\in d}\bigl[(|w_j|^2+8/5)^2-(8/5)^2\bigr]
 +\sum_{j\in W}\bigl[(|w_j|^2+9/10)^2-(9/10)^2\bigr],
 \label{eq:v7-relative-norm-edge-action}
\end{equation}
где $(m_a)=(3,3,1,1,1,1,1)$. Проверяем одну положительную сумму
\begin{equation}
 \mathcal S_\beta=\mathcal S_E+\beta\mathcal S_V,
 \qquad \beta\ge0.
 \label{eq:v7-relative-norm-combined-action}
\end{equation}

\section{Запуск нулевого сектора}

Полный гессиан в нуле вычислен на той же вещественной размерности
$7+20=27$. Условие положительности всех тяжёлых направлений имеет точную
границу
\begin{equation}
 0\le\beta<\frac8{15}.
 \label{eq:v7-relative-norm-beta-window}
\end{equation}
При естественном кандидате $\beta=1/2$ получаем
\begin{equation}
 (n_-,n_0,n_+)_{0}=(7,0,20),
 \qquad \lambda_{\rm heavy}^{\min}=0.4.
 \label{eq:v7-relative-norm-origin-signature}
\end{equation}
Ровно семь корневых направлений запускаются, а все двадцать down--weak
конкурентов получают положительную щель. Равный вес $\beta=1$ этот тест не
проходит.

\section{Целевой вакуум}

В точке $A=A_0$, $|z_a|=1$, $w_j=0$ обе положительные нормы минимальны.
Полный гессиан равен
\begin{equation}
 (n_-,n_0,n_+)_{*}=(0,0,27),
 \qquad \lambda_{\min}=5.6.
 \label{eq:v7-relative-norm-vacuum-signature}
\end{equation}
Таким образом, положительная кривизностная архитектура одновременно даёт
селективный запуск и строго устойчивый локальный вакуум без отрицательной
метрики или теплового суперследа.

\begin{theorem}[Одноотношенческий положительный локальный проход]
На физическом incidence-блоке $10\times11$ и выведенном рёберном Hodge-блоке
действие \eqref{eq:v7-relative-norm-combined-action} при $\beta=1/2$ имеет
правильную нулевую сигнатуру $(7,0,20)$ и положительную вакуумную сигнатуру
$(0,0,27)$. Допустимое окно нормировки открыто и заканчивается точно при
$8/15$.
\end{theorem}

\section{Оставшийся нормировочный разрыв}

Число $1/2$ естественно напоминает физический полуслед Real-пары и
стандартный коэффициент квадрата кривизны. Однако оно ещё не выведено из
одного представленного модуля. Литература по Yang--Mills функционалам
суперсвязностей подтверждает, что квадрат кривизны способен объединять
калибровочный и хиггсовский секторы, но относительные нормировки зависят от
представления и допустимых центральных весов
\cite{v7-relative-norm-roepstorff,v7-relative-norm-figueroa}.

\begin{nogocriterion}[Запрет объявлять $\beta=1/2$ выведенным]
Пока Real-удвоение, физический полуслед и Hodge-метрика обоих блоков не
записаны в одном carrier, $\beta=1/2$ остаётся единственным невыведенным
отношением. Положительный гессиан является условным локальным проходом, а не
полным физическим замыканием.
\end{nogocriterion}

\begin{protocol}
\begin{enumerate}
 \item действие ограничено снизу: \textbf{да};
 \item отрицательная метрика или суперслед: \textbf{нет};
 \item допустимое окно: \textbf{$0\le\beta<8/15$};
 \item кандидат $\beta=1/2$: \textbf{проходит};
 \item нулевая сигнатура: \textbf{$(7,0,20)$};
 \item тяжёлая щель в нуле: \textbf{$0.4$};
 \item вакуумная сигнатура: \textbf{$(0,0,27)$};
 \item вакуумная минимальная мода: \textbf{$5.6$};
 \item происхождение $\beta=1/2$: \textbf{открыто};
 \item итог: \textbf{одноотношенческий положительный локальный проход}.
\end{enumerate}
\end{protocol}

Машинный сертификат сохранён в
\path{s2t_v7_derived_relative_involution_curvature_norm_gate_results.json}.

\begin{thebibliography}{9}
\bibitem{v7-relative-norm-roepstorff}
G.~Roepstorff,
\emph{Superconnections and the Higgs Field}, arXiv:hep-th/9801040.
\bibitem{v7-relative-norm-figueroa}
H.~Figueroa, J.~M.~Gracia-Bondia, F.~Lizzi, J.~C.~Varilly,
\emph{A Nonperturbative Form of the Spectral Action Principle in
Noncommutative Geometry}, J. Geom. Phys. 26 (1998), 329--339;
arXiv:hep-th/9701179.
\end{thebibliography}